\chapter{Análisis de impacto}
\label{ch:impacto}

De los contextos en los que un trabajo puede tener repercusión (personal,
empresarial, social, económico, medioambiental y cultural), este se concentra
por su naturaleza en dos: el medioambiental, por la especie y el problema que
aborda, y el científico-metodológico, por las contribuciones que deja para
futuros estudios de movimiento animal. Los demás se tratan de pasada. Se cierra
el capítulo valorando la alineación del trabajo con los Objetivos de Desarrollo
Sostenible de la Agenda 2030.

La capacidad de anticipar dónde estará una población de un migrante de larga
distancia tiene un interés claro para la conservación: la planificación de
espacios protegidos, la evaluación del riesgo de colisión con infraestructuras
como los parques eólicos situados en los corredores migratorios o el seguimiento
de las rutas son tareas que se benefician de modelos predictivos. Conviene, sin
embargo, ser realista sobre el alcance de este trabajo concreto. Su aporte
medioambiental no reside en una predicción precisa del destino, sino en señalar los
días en que el ave está migrando, acompañando la predicción de una banda de
incertidumbre calibrada, y en ofrecer una herramienta (capítulo \ref{ch:o5}) que
permite visualizar ese comportamiento sobre el mapa.

El impacto más tangible del trabajo es, no obstante, metodológico, y se proyecta
sobre cualquier estudio que prediga movimiento animal a partir de rastros GPS.
Tres aportaciones son reutilizables más allá de este caso. La primera es el
diseño causal de las variables: exigir que toda variable predictiva sea
calculable sin observar el futuro evita una fuga de información que, de otro modo,
inflaría de forma artificial las métricas, y es una disciplina trasladable a
cualquier problema de predicción sobre series temporales. La segunda es el uso de
la coherencia biológica como criterio de validación, contrastando todo resultado
contra la fenología conocida de la especie antes de darlo por bueno. La tercera es
la caracterización explícita de un techo estructural: mostrar con datos que el límite
de la predicción a un día no proviene de la configuración de los modelos, sino del
horizonte y de la información disponible, es en sí mismo un resultado que orienta
dónde merece la pena invertir el esfuerzo futuro (modelos de secuencia, viento o
predicción multi-paso). A ello se suma que todo el trabajo se ha construido de
forma reproducible, con cada decisión documentada y respaldada por su evidencia.

En lo personal, el trabajo ha consolidado la formación del autor en modelado
probabilístico y aprendizaje automático aplicados a un problema real. En lo
económico y empresarial su repercusión directa es escasa, propia de un estudio de
carácter académico, si bien el tipo de herramienta desarrollada encaja con las
necesidades de organismos de conservación y consultoras ambientales.

Respecto a los Objetivos de Desarrollo Sostenible de la Agenda 2030, el trabajo se
alinea sobre todo con el objetivo 15 (vida de ecosistemas terrestres), cuya meta de
proteger la biodiversidad y frenar su pérdida se sirve de herramientas que ayuden a
comprender y conservar las especies migratorias. En segundo lugar conecta con el
objetivo 13 (acción por el clima): el calendario y las rutas de la migración son
sensibles a las condiciones ambientales, de modo que los modelos que caracterizan
el movimiento pueden servir, a más largo plazo, como indicadores de los
desplazamientos fenológicos asociados al cambio climático. De forma más tangencial,
la herramienta de cartografía interactiva tiene un valor divulgativo que la acerca
al objetivo 4 (educación de calidad).
